\chapter*{Abstract}
\addcontentsline{toc}{chapter}{Abstract}

SavePoint grew out of applying to video games the personal cataloguing pattern of platforms such as Goodreads or Letterboxd, adding an inventory of physical and digital ownership that those products never had to solve; on that basis, the project turns the platform into a reproducible testbed for recommender systems. The web application offers a catalogue of more than one hundred and ninety thousand governed works sourced from the Internet Game Database, a personal library, physical and digital copies, lists, comments, profiles with privacy settings, friendships and recommendations, with an accessible interface in Spanish and English that adapts to desktop and mobile. It is verified with more than eight hundred automated tests and has been deployed publicly on free services, although with a subset of the catalogue. The application and the reproducible evaluation of the recommenders are the two complementary halves of the work: the former collects and governs the data the latter needs.

The system compares sixteen recommendation algorithms: two reference baselines, content-based variants, neighbourhood collaborative filtering, and hybrid approaches with diversity re-ranking. This report describes their formulas, item and user representations, and hybrid combination. The offline protocol fixes corpus, versioned snapshots, seeds, shared candidates, exclusions and ranking metrics at $K\in\{5,10,20\}$.

The published evidence, version 15 of the protocol, uses a population of 400 simulated users, 79 of them evaluable in the test partition, with a single test run. In it, content-based and hybrid variants outperform the reference algorithms with statistical significance: the best content-based variant returns at least one of the held-out games within the first ten positions for 49\,\% of the evaluable users, against 3\,\% for popularity recommendation. The evidence supports an internal comparison under the frozen corpus and protocol, but it does not represent real users or justify automatic generalisation. This report explicitly separates that publication from the later implementation contract.

An agent-assisted engineering methodology governed through the Get Stuff Done method and conceptually organised with Obsidian was applied throughout the project. The tools were used under human direction and review, while final decisions and responsibility for analysis, design, implementation and validation remained with Felipe Peña Núñez.

\vspace{.5cm}
\textbf{Keywords:} video games, recommender systems, offline evaluation, reproducibility, collaborative filtering, content-based recommendation, data provenance.
